% auto-generated by scripts/build_cost_sensitivity_from_profit_long.py
\begin{table}[htbp]
  \centering
  \small
  \caption{コストシナリオ別の最頻推奨処置（\texttt{profit\_long.csv} の $p$・\texttt{history} を固定し、係数のみ変更して再計算。参考）}
  \label{tab:cost_sensitivity_modes}
  \begin{tabular}{@{}p{0.42\linewidth}p{0.48\linewidth}r@{}}
    \toprule
    シナリオ（概略） & 最頻推奨（処置） & 割合 \\
    \midrule
    extreme\_lean（割引係数小） & \texttt{Discount / Web} & 86.1\% \\
    low\_cost（config 相当・やや控えめ） & \texttt{No Offer / Web} & 96.3\% \\
    mid\_cost & \texttt{No Offer / Web} & 100.0\% \\
    high\_cost & \texttt{No Offer / Web} & 100.0\% \\
    extreme\_heavy（割引係数大） & \texttt{No Offer / Web} & 100.0\% \\
    \bottomrule
  \end{tabular}
\medskip
  \par\noindent\footnotesize 極端設定間（\texttt{extreme\_lean} vs \texttt{extreme\_heavy}）で\textbf{同一の離散処置}が選ばれた顧客の割合：\textbf{13.9\%}（残りは処置が不一致＝コスト仮定で結論が入れ替わる）。
\end{table}
